% Introduction to Part X. Unnumbered text: no \label except the \phantomsection target.
\phantomsection\label{intro:X}
Part~X puts one small-field step into polymer form (Sections~\ref{sec:complex-references}--\ref{sec:cluster}, Theorem~\ref{thm:L}) and iterates it over all scales (Sections~\ref{sec:flow}--\ref{sec:flow4}, Theorem~\ref{thm:N}); Remark~\ref{rem:flow-status} describes the organization. The results concern the modified small-field model on a fixed torus. At every step the integration is restricted to small fields and the rest is discarded, so the partition function of this model is at most that of the lattice theory with the same counterterm $E_0$ (Remark~\ref{rem:sf-reference}). The large-field stage is not treated.

\paragraph{The step in polymer form.} The activities must be analytic in the background on complex domains whose constants do not depend on the coupling, the step, the torus or the polymer (Remark~\ref{rem:domain-constants}).
\begin{itemize}
\item Section~\ref{sec:complex-references} extends Parts~IV and~IX to complex references (Propositions~\ref{prop:AB} and~\ref{prop:complex-ext}), and Section~\ref{sec:step-quantities} turns the ingredients of Theorem~\ref{thm:step} into functions of the weakening parameters that factorize over polymers (Propositions~\ref{prop:interaction-pieces}--\ref{prop:hessian-general}), and localizes them by Lemma~\ref{lem:localization} (\hyperref[intro:IX]{introduction to Part~IX}).
\item Section~\ref{sec:cluster} imposes the cutoffs on the fluctuation field and on the white noise (Definition~\ref{def:small-field-step}) and evaluates the fluctuation integral by a Mayer and polymer expansion (Theorem~\ref{thm:mayer}). By parity in the coupling there is no term of order $g$, and the pure sources are of order $g^2p_0^6$ (Proposition~\ref{prop:pure-sources}). The result is Theorem~\ref{thm:sf-step}.
\end{itemize}

\paragraph{The flow.} In four dimensions the coupling is marginal, and the activities are of order one, not small in the coupling (Remark~\ref{rem:flow-status}). So the old activities are not reblocked and contracted; the density keeps Ba\l aban's generation form \cite[(0.22)--(0.30)]{BalYM4I}.
\begin{itemize}
\item A generation of scale $j$ is a family of normalized, gauge-invariant activities on the polymers of the grid of scale $j$, bounded only at that scale (Definition~\ref{def:generation}). The new terms of a step form a raw generation (Lemma~\ref{lem:raw-new}). Its average over the symmetries of its grid (Lemma~\ref{lem:symmetrization}) has an exactly isotropic period density, $\beta_E$ times that of the action (Proposition~\ref{prop:period-iso}), and $g_{k+1}^{-2}=g_k^{-2}-\beta_{k+1}$.
\item At every later scale the marginal part of an old generation is extracted: after $m\ge m_0$ scales, its small polymers in a grid cube sum to $\beta_jS_\Box$ up to $C\mathsf KL^{-\alpha m}$ (Proposition~\ref{prop:extraction}; $\mathsf K$ is the bound on the generations, Remark~\ref{rem:flow-constants}), and its other polymers carry a factor $L^{-m}$ (Lemma~\ref{lem:large-later}). Grouped with their counterterms $-\beta_jS$, the old generations form old activities bounded uniformly in the scale (Proposition~\ref{prop:old-activities}), and no contraction is needed.
\item Sections~\ref{sec:flow} and~\ref{sec:target-flow} fix the parameters (Remark~\ref{rem:parameters}), the grids, the torus, the first step and the vacuum energy (Propositions~\ref{prop:first-step} and~\ref{prop:vacuum}).
\end{itemize}
One step preserves the generation form (Theorem~\ref{thm:gen-step}), and Theorem~\ref{thm:sf-flow} iterates it. As long as the couplings stay in $(0,\gamma]$: the density of every scale is in generation form, every generation is bounded at its creation scale, $|\beta_j|\le C_\beta\mathsf K$ and $\frac34g_{j-1}\le g_j\le2g_{j-1}$, with $\mathsf K$ and $C_\beta$ independent of the scale, the torus and the couplings; and, for the counterterm $E_0$ with $\eps_K=0$, the partition function of the small-field model is the explicit product $\mathcal Z_K$ of normalizations times a number between $e^{-c|\T^{(K)}|}$ and $e^{c|\T^{(K)}|}$, with $c$ depending on $g_K$ but not on $K$, the torus or the earlier couplings (Corollary~\ref{cor:sf-stability}). The generations and the $\beta_j$ depend continuously on the couplings (Proposition~\ref{prop:flow-continuity}). The theorem assumes no asymptotic behaviour of the couplings and gives no sign of the $\beta_j$ (Remark~\ref{rem:sf-flow}(a); Theorem~\ref{thm:sf-flow}(c) bounds only $|\beta_j|$); Part~XI derives the behaviour of the couplings from the one-loop coefficient (Theorem~\ref{thm:M}(iv)).

\paragraph{Imports and exports.} Part~X uses Theorem~\ref{thm:step}, Lemmas~\ref{lem:eliminate} and~\ref{lem:gf-step-invariance} (Part~VII), the graded domains and Lemma~\ref{lem:quadratic-bound} (Part~VIII), Corollary~\ref{cor:weak-covariance}, Proposition~\ref{prop:one-loop-polymer} and Theorem~\ref{thm:weak-bg} (Part~IX), Dimock's scheme \cite{DimockI} and the Koteck\'y--Preiss criterion \cite{KP}. The flow goes to Part~XI, which proves its properties (F1)--(F3) (Proposition~\ref{prop:flow-F}) and the ultraviolet behaviour of its couplings (Theorem~\ref{thm:ol-uv}). Part~XI also uses Lemma~\ref{lem:first-step-formula}, Proposition~\ref{prop:period-iso}, and Proposition~\ref{prop:jacobians}(d) and condition~($\Gamma$4) of Remark~\ref{rem:flow-constants}, which serve only Part~XI. Remark~\ref{rem:symmetrization-flow}(b) is not needed for the flow.

\paragraph{Details for this Part.} The companion has two chapters on Part~X. The first expands the complex references, the step quantities, the cluster expansion and the reblocking; the second the parameters and the grids, the generations, the period density, the extraction, the first step and the flow.
